%% =================================================================
%%  3.  The Unified Algebraic-Lattice Bipartition Framework
%% =================================================================
\section{The Unified Algebraic-Lattice Bipartition Framework}
\label{sec:ualbf}

By \cref{thm:qpn-odd-sq}, every QPN is an odd perfect square
$N = \prod_{i} p_i^{2e_i}$
with each $p_i$ an odd prime.  The key idea of the UALBF is to
\emph{bipartition} this factorisation into a \emph{prefix} $N_L$ and a
\emph{suffix} $N_R$, chosen so that all arithmetic constraints induced by the
QPN equation decompose cleanly across the split.  The prefix is constructed
incrementally by the Rust engine's depth-first search, while the suffix's
residue class modulo $\sigma(N_L)$ is determined by a constant-time
modular inverse; candidates in that class are then enumerated and sieved.

%% -------------------------------------------------------------------
\subsection{Bipartitioning the Search Space}
\label{subsec:bipartition}

\begin{definition}[QPN Bipartition]\label{def:bipartition}
A \textbf{QPN bipartition} of a quasiperfect number $N$ is a triple
$(N, N_L, N_R)$ satisfying:
\begin{enumerate}
  \item $N_L > 0$, $N_R > 0$;
  \item $N = N_L \cdot N_R$;
  \item $\gcd(N_L, N_R) = 1$.
\end{enumerate}
Any partition of the prime set of $N$ into two disjoint subsets induces
a valid bipartition; the Rust engine's DFS assigns the first $k$ primes
to $N_L$ and all remaining primes to $N_R$.
\end{definition}

\paragraph{Multiplicativity over the bipartition.}
Because $\gcd(N_L, N_R) = 1$, the sum-of-divisors function factors:
\begin{equation}\label{eq:sigma-bipartition}
  \sigma(N) \;=\; \sigma(N_L)\,\sigma(N_R).
\end{equation}
Substituting the QPN equation $\sigma(N) = 2N + 1 = 2 N_L N_R + 1$ yields
the \emph{fundamental bipartition identity}:
\begin{equation}\label{eq:bipartition-identity}
  \sigma(N_L)\,\sigma(N_R) \;=\; 2\,N_L\,N_R + 1.
\end{equation}

\paragraph{Coprimality of prefix and its divisor sum.}
The identity~\eqref{eq:bipartition-identity} has an important structural
consequence:

\begin{theorem}[Prefix--Sigma Coprimality]\label{thm:prefix-coprime}
For any QPN bipartition $(N, N_L, N_R)$,
\[
  \gcd\bigl(N_L,\;\sigma(N_L)\bigr) \;=\; 1.
\]
\end{theorem}

\begin{proof}
Let $d = \gcd(N_L, \sigma(N_L))$.  We show $d \mid 1$.

Since $d \mid N_L$, we have $d \mid N_L \cdot N_R$, hence
$d \mid 2\,N_L\,N_R$.
Since $d \mid \sigma(N_L)$ and $\sigma(N_L)$ is a factor of
$\sigma(N_L)\cdot\sigma(N_R)$, we have $d \mid \sigma(N_L)\cdot\sigma(N_R)$.
By the bipartition identity~\eqref{eq:bipartition-identity},
$\sigma(N_L)\cdot\sigma(N_R) = 2N_LN_R + 1$.
So $d$ divides both $\sigma(N_L)\cdot\sigma(N_R)$ and
$2N_LN_R$, hence $d$ divides their difference:
$d \mid (2N_LN_R+1) - 2N_LN_R = 1$.
Therefore $d = 1$.
\end{proof}

\cref{thm:prefix-coprime} is operationally critical: it guarantees that the
modular inverse of $2 N_L$ modulo $\sigma(N_L)$ always exists, so the Rust
engine's ray-cast subroutine can never encounter a division-by-zero panic.

\paragraph{The AMBS suffix target.}
Casting~\eqref{eq:bipartition-identity} into the ring
$\mathbb{Z}/\sigma(N_L)\mathbb{Z}$, the left-hand side vanishes (since
$\sigma(N_L) \equiv 0$), giving
\[
  0 \;=\; 2\,N_L\,N_R + 1
  \qquad\text{in } \mathbb{Z}/\sigma(N_L)\mathbb{Z}.
\]
Rearranging:
\begin{equation}\label{eq:ambs-target}
  N_R \cdot (2\,N_L) \;\equiv\; -1 \pmod{\sigma(N_L)}.
\end{equation}
Because $\gcd(2 N_L, \sigma(N_L)) = 1$ (by \cref{thm:prefix-coprime} and
the fact that $\sigma(N_L)$ is odd for a QPN), the value of $N_R$ modulo
$\sigma(N_L)$ is uniquely determined.  This is the \emph{Algebraic-Modular
Bipartition Sieve} (AMBS) target: given a prefix $N_L$, the engine computes
the unique residue class of $N_R$ in $\mathcal{O}(1)$ time via a single
modular inverse, then enumerates all representatives
$z \equiv r \pmod{\sigma(N_L)}$ in the interval $[z_{\min},\, z_{\max}]$
in $\mathcal{O}(z_{\max}/\sigma(N_L))$ time, applying the mod-8 sieve
and coprimality checks to each candidate.

The AMBS target equation~\eqref{eq:ambs-target} also admits a clean
contrapositive:

\begin{theorem}[Ray-Cast Exhaustion Certificate]\label{thm:no-sol-no-qpn}
If $N_R \cdot (2\,N_L) \not\equiv -1 \pmod{\sigma(N_L)}$ for all
valid suffix candidates, then $N = N_L \cdot N_R$ is not quasiperfect.
\end{theorem}

\begin{proof}
Suppose for contradiction $N$ is quasiperfect.
By~\eqref{eq:ambs-target}, any such $N$ must satisfy
$N_R \cdot (2N_L) \equiv -1 \pmod{\sigma(N_L)}$.
But every valid suffix candidate fails this congruence by hypothesis,
so the QPN hypothesis implies a contradiction.
\end{proof}

\cref{thm:no-sol-no-qpn} is the formal certificate that the Rust engine's
ray-cast exhaustion constitutes a {\em complete} non-existence proof:
when the engine enumerates all representatives in the residue class
$z \equiv r \pmod{\sigma(N_L)}$ within the search bounds and none
satisfies all divisibility constraints, the prefix $N_L$ is formally
proven to be incompatible with any QPN\@.  In the Lean formalisation,
this is
\texttt{no\_\allowbreak{}solution\_\allowbreak{}no\_\allowbreak{}qpn}
in \texttt{Bipartition.lean}---the mathematical proof that the Rust
ray-cast is a complete, rigorous exhaustion of the search space, not
merely a heuristic search.

%% -------------------------------------------------------------------
\subsection{Parity Contradictions for Even Squares}
\label{subsec:double_squares}

The restriction to odd perfect squares established in \cref{thm:qpn-odd-sq}
rests on two independent contradictions.  The modulo-3 argument
(\texttt{qpn\_\allowbreak{}not\_\allowbreak{}double\_\allowbreak{}square}, cf.\ \cref{subsec:sigma_properties})
eliminates the double-square form $N = 2m^2$ outright.  Here we detail the
deeper Legendre-symbol obstruction that eliminates even perfect squares.

\begin{theorem}[Even Square QPN Obstruction]\label{thm:even-sq-obstruction}
If $N = m^2$ is a QPN and $m$ is even, then a contradiction arises.
\end{theorem}

\begin{proof}
Write $m = 2^e u$ with $u$ odd and $e \ge 1$, so
$N = 2^{2e}\,u^2$.  The QPN equation and multiplicativity of $\sigma$ give
\[
  (2^{2e+1}-1)\,\sigma(u^2) \;=\; 2^{2e+1}\,u^2 + 1.
\]
Set $M = 2^{2e+1}-1$.  Rewriting the right-hand side as
$(M+1)\,u^2 + 1 = M\,u^2 + u^2 + 1$, we obtain
\begin{equation}\label{eq:M-divides}
  M \mid (u^2 + 1).
\end{equation}
Since $2e+1 \ge 3$, $2^{2e+1} \equiv 0\pmod 4$, so $M = 2^{2e+1}-1 \equiv 3 \pmod{4}$.

\textbf{Claim.}  Any integer $\equiv 3 \pmod 4$ has a prime factor $q \equiv 3 \pmod 4$.
\emph{Proof of claim:}  by strong induction.  If $M$ is prime, $q = M$.  Otherwise
$M = ab$ with $1 < a,b < M$.  Since $(ab) \bmod 4 = 3$, examining
$a \bmod 4 \times b \bmod 4 \pmod 4$ shows at least one of $a, b$ is $\equiv 3 \pmod 4$
(products $1\times 1 = 1$, $1\times 3 = 3$, $3\times 3 = 1$, etc.), and the
factor $\equiv 3 \pmod 4$ is smaller, so the inductive hypothesis applies.

Applying the claim, there exists a prime $q \mid M$ with $q \equiv 3 \pmod{4}$.
By~\eqref{eq:M-divides} and transitivity, $q \mid u^2 + 1$, so
$u^2 \equiv -1 \pmod{q}$.
But by the first supplement to quadratic reciprocity, $-1$ is a quadratic
residue modulo a prime $q$ if and only if $q \equiv 1 \pmod{4}$.  Since
$q \equiv 3 \pmod{4}$, this is impossible.
\end{proof}


Combined with the modulo-3 elimination of $N = 2m^2$,
\cref{thm:even-sq-obstruction} closes every non-odd-square alternative.
In the Lean formalisation, this overarching result is verified as
\texttt{square\_\allowbreak{}qpn\_\allowbreak{}parity\_\allowbreak{}obstruction}
in \texttt{QPN/BasicProperties.lean}.
The UALBF therefore searches exclusively over factorisations of the form
$N = \prod_i p_i^{2e_i}$ with each $p_i$ odd, with the bipartition
machinery of this section reducing each candidate to a single modular
equation~\eqref{eq:ambs-target}.
